\section{Component and Semiconductor Stress}
\label{sec:analysis-stress}

The half bridge is exposed to DC-link voltage plus transient overshoot. Its current is the signed sum of pump current and the complete primary resonant-path current:
\begin{equation}
i_{\mathrm{HB}}(t)=i_P(t)+i_{\mathrm{res,p}}(t).
\end{equation}
The primary branch already includes magnetizing current: \(i_{\mathrm{res,p}}=i_{\mathrm{load,p}}+i_m\). Adding \(i_m\) to the complete branch again would double-count it.

Adding the branch-current magnitudes gives a conservative peak estimate. Calculating RMS loss requires their actual signed waveform through active and idle intervals. For the triangular pump approximation,
\begin{equation}
I_{P,\mathrm{pk}}\approx4f_s\CP V_{\mathrm{in,pk}}.
\end{equation}
Pump current can impose substantial half-bridge stress at modest connector current. Calibrated primary-current feedback at PA2 supports load classification; it does not directly measure connector current.

Pump-capacitor, inductor, and steering-diode stress depends on line phase, link voltage, and charge-transfer completion. Isolated-branch circulating current contributes winding and resonant-capacitor heating, while secondary rectifier stress depends on conduction intervals and storage voltage.

Driving the original prototype's burst-enable output PA10 low inhibits excitation but does not remove stored tank energy. The hard OVP margin must therefore include sensing and actuation delays and residual transfer after inhibition.

Measured bus maxima, overshoot, current peaks, and temperature are compared with the \SI{450}{\volt} capacitor rating and the semiconductor ratings. The initial calculations use nominal passive values and supplied magnetic measurements. The waveforms give the stresses reached in the tested circuit.
